\documentclass[11pt,a4paper]{article}
\usepackage[T1]{fontenc}
\usepackage[utf8]{inputenc}
\usepackage{lmodern}
\usepackage{amsmath,amssymb}
\usepackage[margin=25mm]{geometry}
\usepackage[hidelinks]{hyperref}
\hypersetup{pdftitle={SU(2) on the full mid-plane: the relative-order-t halving defect},pdfauthor={navstokgap research working notes}}
\newcommand{\tightlist}{\setlength{\itemsep}{0pt}\setlength{\parskip}{0pt}}
\setlength{\emergencystretch}{3em}
\setcounter{secnumdepth}{0}
\begin{document}
\hypertarget{su2-on-the-full-mid-plane-the-relative-order-t-halving-defect}{%
\section{SU(2) on the full mid-plane: the relative-order-t halving
defect}\label{su2-on-the-full-mid-plane-the-relative-order-t-halving-defect}}

\textbf{Result, 2026-09-27 (formal one-loop calculation).} For one
halving of direction 1, with isotropic coarse heat time
\(t=\lambda_3a\), the zero-momentum coupling coefficients in Hypothesis
P(\(\alpha\)) are

\[\boxed{\begin{aligned}
\delta_t^{12}=\delta_t^{13}
 &=\frac{t}{24}\int_{\mathcal B}\frac{q}{R}+O(t^2)>0,\\
\delta_t^{23}
 &=t\left(-\frac1{12}+2\mathcal I\right)+O(t^2)<0,
\end{aligned}}\tag{1}\]

where the inequalities describe the leading coefficients, and

\[\begin{gathered}
\int_{\mathcal B}:=\int_{-\pi}^{\pi}\int_{-\pi}^{\pi}
 \frac{dk_2\,dk_3}{(2\pi)^2},\qquad
q=2-2\cos k_2,\quad r=2-2\cos k_3,\quad d=8+r,\quad R=8+r+q,\\
\boxed{\displaystyle
\mathcal I=\int_{\mathcal B}
\left\{\frac{r^3+10r^2+4r-1248+q(8+r)^2}{24(8+r)^2(8+r+q)}
-\frac{q(2\sin k_3)^2(2\sin k_2)^2}
 {4(8+r)(8+r+q)^4}\right\}<0.}
\end{gathered}\tag{2}\]

All integrals are finite. Their denominators are bounded below by
positive powers of 8. The first numerator in (2) increases with each of
\(q,r\in[0,4]\) and has maximum \(-432\); the second summand is
nonpositive. The cut integral is strictly positive. These statements use
no numerical evaluation.

The cut shifts come from the normalized bridge-curvature determinant
(N2). For equal adjacent layers, the bridges have an exactly Gaussian
zero-image density in exponential coordinates. The transverse shift then
consists of the mid-face heat-kernel amplitude, \(-t/12\), and the
coupled noncommutative fluctuation determinant, \(2t\mathcal I\) (N3,
including the curvature of the mid-face action). Neighbouring faces
share fluctuations throughout this calculation.

After these shifts and the vacuum term are extracted, the formal
\textbf{bulk local} smooth-field expansion contains only operators of
geometric dimension at least six. Here \([D]=1\) and \([F]=2\), the
convention used for irrelevance in the
\href{series-parallel-gauge-refinement.md}{series/parallel note}. The
full small-field statement still requires the volume-uniform, normalized
estimates in §7. Equation (1) gives bulk coefficients; finite tori also
carry the explicit flat-holonomy term (16) below. It obstructs the
literal flux-only P(\(\alpha\)) inequality at fixed torus size as
\(t\to0\), and is exponentially small when the number of sites around
the torus tends to infinity. Iteration requires a further stability
estimate.

\textbf{Inputs and scope.} The exact definition of \(\Psi\) and
\(\mathcal D\) is Proposition 1, equations (3)--(4), and Hypothesis
P(\(\alpha\)) of the
\href{series-parallel-gauge-refinement.md}{series/parallel note}
(full-read: relevant sections, including Propositions 4, 6 and 7). The
bridge checks use \href{su2-midpoint-exact.md}{the exact midpoint note},
Theorems 1--4 (full-read). The calculations below are written
derivations of the formal coefficients; uniformity in the spin sum is a
separate estimate (§7). We use \([T_a,T_b]=\epsilon_{abc}T_c\),
\(|T_a|=1\), \(C_2(\mathrm{adj})=2\), and normalized Haar measure.

\hypertarget{expansion-and-the-massive-mid-plane-covariance}{%
\subsection{1. Expansion and the massive mid-plane
covariance}\label{expansion-and-the-massive-mid-plane-covariance}}

Use the mid-vertex gauge of Proposition 1 and write
\(m_e=m_{*e}\exp\xi_e\). At the trivial background the quadratic
exponent is

\[\frac1t\left(2\sum_e|\xi_e|^2
 +\frac14\sum_g|\bar\Phi_g+(C\xi)_g|^2\right),\qquad
\bar\Phi=\frac{\Phi_\ell+\Phi_{\ell+1}}2.\]

The Fourier symbol of the curl and the inverse fluctuation Hessian are

\[c(k)=(1-e^{ik_3},\ e^{ik_2}-1),\qquad
H_0=4I+\frac12c^*c,\qquad
G=H_0^{-1}=\frac14\left(I-\frac{c^*c}{R}\right).\tag{3}\]

There is an additional identity matrix on the three colour components.
In particular \(4I\le H_0\le8I\), including zero momentum. The bridge
term supplies a mass for every mid-edge fluctuation. Integrating (3)
gives precisely the free defect \(\mathcal D_0\) of Proposition 2; the
plane coupling agrees with the covariance in Proposition 4.

For precision about orders, introduce local boundary coordinates \(X\)
and the zero-image exponent \(\mathcal S(X,\xi)\) of the
\emph{normalized} integral defining the defect. Thus it includes the
subtraction \(-\sum_e|X_e|^2/2\) from the bridge denominators and
\(-\sum_{f\perp1}|X_f|^2/4\) from the old-face ratio. Its positive terms
are the two squared half-face distances per cut face and
\(\sum_g|\log m_{\partial g}|^2/4\). Let
\(F(X)=\min_\xi\mathcal S(X,\xi)\). The formal normalized Laplace
expansion of the bulk part is

\[\mathcal D(X,t)=c_{\rm vac}(t)N+\frac{F_2+F_3+F_4}{t}
 +A_2(X)+\text{higher terms},\qquad F_2/t=\mathcal D_0.\tag{4}\]

Finite periodic planes also have the winding functional of §7.
Subscripts are homogeneous boundary degree about the trivial field. The
one-loop amplitude \(A\) includes \(\tfrac12\log\det H\), Haar and
heat-kernel amplitudes, and the bridge denominators. An explicit
prescription, useful also for its derivative terms, is

\[\begin{aligned}
A(X)={}&\tfrac12\log\det H(X)
 +\tfrac12\sum_e[\log j(L_e^-)+\log j(L_e^+)-\log j(X_e)]\\
 &-\sum_e\log j(\xi_e^*)+\tfrac12\sum_g\log j(Y_g^*),
\qquad j(Z)=\left[\frac{\sin(|Z|/2)}{|Z|/2}\right]^2.
\end{aligned}\]

Here \(\xi^*\) is the classical saddle, \(L_e^\pm\) are the logarithms
of its two half-face holonomies, \(Y_g^*\) is its mid-face logarithm,
and \(H\) is the Hessian in the \(\xi\) coordinates.
Boundary-independent constants are suppressed. Gauge invariance removes
the linear term of the bulk local expansion. Its quadratic term is

\[A_2(X)=\frac1{48}\int_{\mathcal B}\frac qR
 \sum_{f\in(12),(13)}|X_f|^2
 +\left(-\frac1{24}+\mathcal I\right)
 \sum_{f\in(23)}|X_f|^2+A_{2,\mathrm{der}}(X).\tag{5}\]

The last quadratic kernel has zero constant term in its covariant
external-momentum expansion. Equation (5) defines the convention for
coupling shifts: the coefficient of \(|X_f|^2\) is
\(\delta_t^{\mu\nu}/(2t)\). On a finite periodic plane, momentum
integrals for local contractions become sums; (1)--(2) specify their
bulk limit.

For typical perturbative counting \(X=O(\sqrt t)\), \(F_3/t\) has order
\(\sqrt t\), and both \(F_4/t\) and \(A_2\) have order \(t\). Keeping
the cubic and quartic classical terms is essential to an order-\(t\)
expansion. On P(\(\alpha\))'s larger set, \(|X|\le t^{1/2-\delta}\),
retain their actual powers of \(X\) rather than replacing \(X^2\) by
\(t\).

\hypertarget{n2-curvature-of-the-normalized-bridges}{%
\subsection{2. N2: curvature of the normalized
bridges}\label{n2-curvature-of-the-normalized-bridges}}

Proposition 6 in its group-general form gives

\[h_X=I+\frac{(\operatorname{ad}X)^2}{48}+O(|X|^4),\qquad
\operatorname{Cov}\xi=\frac t4h_X^{-1}+O_X(t^2).\tag{6}\]

For \(SU(2)\), \(h_X\) has eigenvalues \(1,h,h\), with
\(h=(|X|/4)\cot(|X|/4)\). The trace of its covariance excess is
\(t|X|^2/96\), or \(t|X|^2/288\) per colour direction. This is the
softening in Proposition 6, with its axial part retained.

Theorems 1--2 of the midpoint note check (6) in every fixed-spin
character. Theorem 3 gives the scalar spin-\(1/2\) matrix; Theorem 4
gives distinct spin-1 entries and their image bounds. At \(X=0\) there
is a stronger useful identity. In exponential coordinates the zero-image
factors satisfy

\[k_{t/2}(e^\xi)^2\,dm
 \ \propto\ e^{-2|\xi|^2/t}\,d^3\xi.\tag{7}\]

Indeed, the two factors \((|\xi|/2)/\sin(|\xi|/2)\) cancel the Haar
Jacobian \([\sin(|\xi|/2)/(|\xi|/2)]^2\). Continuing this Gaussian to
\(\mathbb R^3\) gives covariance exactly \(tI/4\); truncation to the
injectivity chart changes it by an exponentially small tail. Its
fundamental expectation \(e^{-t/32}(1-t/16)\) and spin-1 expectation
\([1+2e^{-t/8}(1-t/4)]/3\) agree with the exact endpoint limits. In
these coordinates the central-bridge covariance is exactly \(tI/4\) up
to the exponentially small tail. (Convention: §1 writes \(m_e=m_*e^\xi\)
and §3 transports as \(m_e=e^\xi U_e\); the two lists differ by a
diagonal unitary, which leaves every determinant unchanged. Abelian
limit: \((\operatorname{ad}X)^2=0\), no commutator and no Jacobian
factor, so all three shifts vanish, matching the free-field check of
P(\(\alpha\)) and Propositions 2 and 4 of the series/parallel note.
Fable referee, 2026-09-27.)

To extract the \(12\) coefficient, take constant commuting cut flux
\(X_{12}=X\), \(X_{13}=0\), and flat midpoint connection. Adjacent
layers can be represented by constant links \(e^{-X/2}\) and \(e^{X/2}\)
in direction 2. The classical defect vanishes. The Hessian in (3)
changes by \(4(h_X-I)\) on direction-2 edges, while bridge normalization
supplies \(-\tfrac12\operatorname{Tr}\log h_X\). Therefore its quadratic
amplitude is

\[\begin{aligned}
A_{2,\mathrm{cut}}
 &=\frac1{96}\operatorname{Tr}
 [(4G-I)(\operatorname{ad}X)^2]\\
 &=\frac{|X|^2}{48}\int_{\mathcal B}\frac rR
 =\frac{|X|^2}{48}\int_{\mathcal B}\frac qR.
\end{aligned}\tag{8}\]

Here \(4G_{22}-1=-r/R\) and
\(\operatorname{tr}_{\rm colour}(\operatorname{ad}X)^2=-2|X|^2\). The
exchange of the two integration variables gives the last equality.
Direction 3 gives the same result. Normalization explains the positive
sign: the unconstrained bridge's curvature determinant is subtracted,
and the mid-face weights constrain its broadened directions.

N2 also contributes higher boundary-degree terms. For example, at the
Gaussian saddle \(\xi^{(0)}\) the quadratic bridge-curvature insertion
in the classical exponent is

\[\frac1{24t}\sum_e
 \langle\xi_e^{(0)},(\operatorname{ad}X_e)^2\xi_e^{(0)}\rangle
 =-\frac1{24t}\sum_e|[X_e,\xi_e^{(0)}]|^2.\tag{9}\]

It belongs to \(F_4/t\); its coefficient follows directly from (6).
Other fourth-degree classical terms are retained by the saddle formula
in §6. At equal adjacent layers \(X_{12}=X_{13}=0\), (7) removes the
bridge-shape contribution altogether, which isolates the transverse
calculation below.

\hypertarget{n3-the-full-plaquette-hessian}{%
\subsection{3. N3: the full plaquette
Hessian}\label{n3-the-full-plaquette-hessian}}

Set both adjacent layers equal to a constant Abelian magnetic
background, with \(U_{23}=e^{BT_3}\) and zero cut flux. It is sufficient
to extract the local bulk coefficient at \(B=0\); one can use a large
box before taking this coefficient. The midpoint connection is an exact
stationary point, since neighbouring plaquettes have the same
covariantly constant flux. Its bridge coordinates have density (7).

For four oriented fluctuations transported to a plaquette base, write
\(\eta_1,\ldots,\eta_4\), and \(Y=BT_3\). The quadratic part of
\(L=\tfrac12|\log m_{\partial g}|^2\) is

\[L^{(2)}=\frac12\left\langle\sum_i\eta_i,
 Q_Y\sum_i\eta_i\right\rangle
 +\frac12\sum_{i<j}\langle Y,[\eta_i,\eta_j]\rangle,
\qquad Q_Y=\frac{\operatorname{ad}Y}{2}
 \coth\frac{\operatorname{ad}Y}{2}.\tag{10}\]

This follows by inserting
\(\log(\prod_i e^{\eta_i})=\sum_i\eta_i+\tfrac12\sum_{i<j}[\eta_i,\eta_j] +O(\eta^3)\)
into the squared-distance function. On the charged colour plane
\(Q_Y=h_B I\), \(h_B=(B/2)\cot(B/2)\). The neutral colour has the
unchanged free Hessian. Both the radial Hessian and the ordered
commutator in (10) are needed.

Here is an explicit determinant evaluation retaining both terms. Choose
background links \(U_2=1\), \(U_3(x)=e^{Bx_2T_3}\) and Fourier transform
in direction 3. Combine the two charged real colours into one complex
field \((u,v)\) on edges in directions \((2,3)\). Put

\[s_B=\frac{B/2}{\sin(B/2)},\qquad h_B=s_B\cos(B/2),\qquad
z=k_3+Bx_2,\quad y=z+B/2,\]

and let \(E\) shift \(x_2\) by one. The charged Hessian of \(L\) is the
Hermitian block operator \(K\) with

\[\begin{aligned}
K_{uu}&=2h_B-2s_B\cos y,\\
K_{vv}&=2h_B-s_B(e^{-iB/2}E+e^{iB/2}E^*),\\
K_{uv}&=f(y)E+g(y),\\
f(y)&=s_B(e^{-iB/2}-e^{-iy}),\qquad
 g(y)=s_B(e^{-iy}-e^{iB/2}).
\end{aligned}\tag{11}\]

For example the four transported charged fluctuations are \(u(x)\),
\(v(x+\hat2)\), \(-e^{iB(x_2+1)}u(x+\hat3)\), \(-e^{iB}v(x)\).
Substitution in (10), using \(h_B\pm iB/2=s_Be^{\pm iB/2}\), gives (11).

The fluctuation Hessian in the exponent is \(H=4I+K/2\). The two charged
real colours contribute
\(\tfrac12\operatorname{Tr}_{\mathbb R}\log H =\operatorname{Tr}_{\mathbb C}\log H\).
Thus the desired one-loop coefficient is the \(B^2\) coefficient per
site of \(\log\det(8+K)\). The neutral determinant is constant. This
formula contains the connected two-vertex contribution
\(-\tfrac14\operatorname{Tr}(H_0^{-1}H_1H_0^{-1}H_1)\) in real-colour
notation, together with \(\tfrac12\operatorname{Tr}(H_0^{-1}H_2)\); it
keeps their cancellations in one expression.

\hypertarget{evaluating-the-transverse-coefficient-as-a-lattice-integral}{%
\subsection{4. Evaluating the transverse coefficient as a lattice
integral}\label{evaluating-the-transverse-coefficient-as-a-lattice-integral}}

The \(uu\) block of \(8+K\) is multiplication by

\[D(y)=8+2s_B[\cos(B/2)-\cos y].\]

Eliminating \(u\) gives a scalar tridiagonal Schur complement \(S\). The
elementary identity \(g(y)^*f(y)=-s_Be^{-iB/2}[D(y)-8]\) gives its
forward hopping \(-8s_Be^{-iB/2}/D(y)\). Its diagonal at site \(x_2\) is

\[8+2h_B-s_B^2\left[
\frac{2-2\cos(z+B)}{D(z+B/2)}+
\frac{2-2\cos(z-B)}{D(z-B/2)}\right].\tag{12}\]

A constant unitary phase removes \(e^{-iB/2}\) from the hopping.
Consequently
\(\log\det(8+K)=\operatorname{Tr}\log D+ \operatorname{Tr}\log S\)
exactly at finite volume with compatible boundary conditions.

For the bulk derivative expansion let \(k=k_2\), \(r=2-2\cos z\),
\(q=2-2\cos k\), \(d=8+r\), \(R=d+q\), and \(w=\sin^2z=r(4-r)/4\). The
scalar symmetric, or Weyl, symbol of \(S\) has expansion
\(s_0+B^2s_2+O(B^4)\), where

\[\begin{aligned}
s_0&=\frac{8R}{d},\\
s_2&=-\frac16-\frac{2\cos z}{d}+\frac{4w}{d^2}
 -\frac{2rw}{d^3}-\frac{r^2}{6d^2}-\frac r{6d}
 -\frac{28\cos k}{3d^2}.
\end{aligned}\tag{13}\]

Indeed \(D(z\pm B/2)=d\pm B\sin z-B^2r/12+O(B^3)\), while at a hopping
midpoint \(D(y)=d(y)+B^2(r(y)/24-1/4)+O(B^4)\). These two expansions in
(12) yield (13). The multiplication determinant contributes
\(\int_{\mathcal B}(r-6)/(24d)\) at order \(B^2\).

The trace of the logarithm also includes the symbol-product correction.
Since \(z\) and \(k\) have commutator of magnitude \(B\), the
second-order product of scalar symbols is

\[f\star g=fg+\frac{iB}{2}\{f,g\}
 -\frac{B^2}{8}(f_{zz}g_{kk}-2f_{zk}g_{zk}+f_{kk}g_{zz})+O(B^3).\]

Expanding the resolvent of \(s_0+u\) in this formula and integrating
\(\log s_0=\int_0^\infty[(1+u)^{-1}-(s_0+u)^{-1}]\,du\) gives

\[[B^2]\operatorname{Tr}\log S\,/N
=\int_{\mathcal B}\left[
\frac{s_2}{s_0}+\frac{\det(\partial^2s_0)}{24s_0^2}\right].\tag{14}\]

Here and below \([B^2]\) means the Taylor coefficient. To check the
factor \(1/24\), the resolvent correction is
\(-\det(\partial^2s_0)/(4(s_0+u)^3)+ Q/(4(s_0+u)^4)\), with
\(Q=s_{0,zz}s_{0,k}^2-2s_{0,zk}s_{0,z}s_{0,k} +s_{0,kk}s_{0,z}^2\).
Integration in \(u\) gives \(\det(\partial^2s_0)/(8s_0^2)-Q/(12s_0^3)\).
Integration by parts on the momentum torus gives
\(\int Q/s_0^3=\int\det(\partial^2s_0)/s_0^2\), proving (14).

Two final simplifications expose the sign. Direct substitution gives

\[\frac{r-6}{24d}+\frac{s_2}{s_0}
=\frac{r^3+10r^2+4r-1248+qd^2}{24d^2R}.\]

For \(s_0=8(1+q/d)\), integration by parts first in \(z\) and then in
\(k\) gives

\[\int_{\mathcal B}\frac{\det(\partial^2s_0)}{s_0^2}
=-6\int_{\mathcal B}\frac{q(r')^2(q')^2}{dR^4}.\]

Together these are exactly (2). As an algebraic normalization check,
replacing the 8 in \(8+K\) by a large mass parameter \(M\) gives
\([B^2]\operatorname{Tr}\log(M+K)/N =-1/(3M)+1/(6M^2)+O(M^{-3})\): (11)
yields \([B^2]\operatorname{Tr}K/N=-1/3\) and
\([B^2]\operatorname{Tr}K^2/N=-1/3\) directly.

Finally, the zero-image heat kernel contains
\(j(Y)^{-1/2}=(|Y|/2)/\sin(|Y|/2)\), so the mid-face amplitude
contributes \(-\log j(Y)^{-1/2}=-|Y|^2/24+O(|Y|^4)\) to the action. The
corresponding factors in \(k_{2t}(U_f)/k_t(U_f)\) cancel exactly. Adding
this \(-1/24\) to \(\mathcal I\) and multiplying by \(2t\) proves (1)'s
transverse coefficient in the formal expansion.

\hypertarget{what-the-midpoint-matrices-say-about-n3}{%
\subsection{5. What the midpoint matrices say about
N3}\label{what-the-midpoint-matrices-say-about-n3}}

The exact midpoint results provide a local check and specify their
scope. In spin \(1/2\), the centered expectation is a scalar matrix, so
an isolated cube's mid-face factors into four scalar bridge expectations
times its interpolated holonomy. In spin 1, Theorem 4 gives

\[\lambda_0-\lambda_1
=\frac t8(1-1/h)+O_X(t^2)
=-\frac{t|X|^2}{384}+O(t|X|^4)+O_X(t^2).\]

The unequal eigenvalues retain the axis of each cut flux. Their
commutator is proportional to \([P_i,P_j]\) for transported axial
projectors, and therefore survives for generic oblique axes. This
fixed-spin matrix commutator has order \(t^2\) at fixed side angles, so
the full order-\(t\) plane coefficient needs the plane Hessian of §3.

The heat-kernel character sum samples spins of size \(t^{-1/2}\), and
each shared edge joins neighbouring face representations. The
local-coordinate Hessian (10)--(14) resums the needed quadratic
fluctuations before any fixed-spin truncation. Its linear-in-\(Y\)
commutator vertices have a connected square at quadratic boundary
degree. This is the N3 contribution relevant to a coupling shift.
Theorems 1--4 test the curvature term N2 and the bridge law at \(X=0\)
(\(\lambda(0,t)=1-t/4\)), and the scalar factorization of the isolated
fundamental sector; they give no check on the value of \(\mathcal I\).

For a more detailed split of the determinant, the radial-curvature part
\(Q_Y-I=(\operatorname{ad}Y)^2/12+O(Y^4)\) contributes

\[\mathcal I_{\rm radial}
=-\frac1{12}\int_{\mathcal B}\frac{q+r}{R}.\]

The rest is \(\mathcal I-\mathcal I_{\rm radial}\), including parallel
transport and ordered-commutator vertices and their connected square.
Equations (10)--(14) specify this split unambiguously. The strict
negative sign in (2) concerns the complete determinant. The separate
bridge term N2 is (8); changing fluctuation coordinates can redistribute
intermediate geometric and product terms while leaving (1) fixed.

\hypertarget{the-other-terms-and-their-smooth-field-dimensions}{%
\subsection{6. The other terms and their smooth-field
dimensions}\label{the-other-terms-and-their-smooth-field-dimensions}}

The classical part in (4) can be retained without guessing individual
local counterterms. Expand the explicit squared-distance exponent of §1
as \(\mathcal S_2+\mathcal S_3+\mathcal S_4+\cdots\) and solve its
Gaussian saddle, \(\xi^{(0)}=-G C^*\bar\Phi/2\). Then

\[\begin{aligned}
F_3&=\mathcal S_3(X,\xi^{(0)}),\\
F_4&=\mathcal S_4(X,\xi^{(0)})-
\frac12\langle\partial_\xi\mathcal S_3(X,\xi^{(0)}),
G\,\partial_\xi\mathcal S_3(X,\xi^{(0)})\rangle.
\end{aligned}\tag{15}\]

Thus (4), (5), (10) and (15) define every term through the stated
boundary and loop orders. This includes (9), the classical interpolation
commutator, and the quartic terms from the ordered mid-face product. A
smooth-field example is the N1 term
\(\langle F_{23},[F_{12},F_{13}]\rangle\), of dimension six; quartic
curvature invariants have dimension eight. The free defect and the
quadratic remainder \(A_{2,\mathrm{der}}\) start with covariant
derivatives of curvature, of dimension six.

The reason this list exhausts relevant smooth-field terms is local. The
inverse of (3) has exponential decay in lattice distance, so its formal
connected contractions have analytic external-momentum expansions. Gauge
invariance forces their bulk local terms to be curvature invariants.
Reflection symmetry forbids the odd-parity Chern--Simons term and
dimension-five terms: a dimension-five term has five index slots, so
some coordinate direction occurs an odd number of times and the
reflection of that direction reverses its sign. At dimension four,
independent reflections leave only \(|F_{12}|^2\), \(|F_{13}|^2\) and
\(|F_{23}|^2\); their coefficients are (1). Differences between opposite
faces and mixed quadratic kernels, allowed by Proposition 7, give higher
derivatives in this expansion. Terms with three curvatures start at
dimension six; for \(SU(2)\) the only cubic is the N1 term, since
\(\langle A,[A,B]\rangle=0\), while for \(SU(3)\) the symmetric tensor
\(d_{abc}\) admits a second dimension-six cubic,
\({\rm tr}(F_{12}\{F_{13},F_{23}\})\). This argument uses smooth fields
and the bulk local expansion; arbitrary small plaquettes can still vary
at lattice momenta of order one.

Consequently the mid-plane extension of Proposition 7 has the expected
\textbf{formal} operator content. Equations (1)--(15) leave no extra
relevant bulk operator at relative order \(t\). They leave a
quantitative uniform-integral problem, specified next.

\hypertarget{a-finite-torus-obstruction-and-the-estimates-needed-for-a-theorem}{%
\subsection{7. A finite-torus obstruction and the estimates needed for a
theorem}\label{a-finite-torus-obstruction-and-the-estimates-needed-for-a-theorem}}

The global qualification has a direct test. On an \(N_2\times N_3\)
periodic plane take identical flat adjacent layers with commuting
holonomies \(e^{\theta_2T_3}\) and \(e^{\theta_3T_3}\). Every coarse
plaquette logarithm vanishes. The charged version of (3) has momenta
\(k_j+\theta_j/N_j\), where \(k_j=2\pi n_j/N_j\). Finite-dimensional
Laplace expansion at its unique zero-action saddle therefore gives

\[\begin{aligned}
\mathcal D(U_\theta,t)-\mathcal D(1,t)
 &=\Omega_N(\theta)+O_N(t),\\
\Omega_N(\theta)
 &=\sum_{k}\log\frac{12-2\cos(k_2+\theta_2/N_2)
                          -2\cos(k_3+\theta_3/N_3)}
                         {12-2\cos k_2-2\cos k_3}.
\end{aligned}\tag{16}\]

The bridge distance term has a unique minimum and positive Hessian for
each such fixed flat background. Rescaling the finitely many normal
coordinates by \(\sqrt t\) gives (16); the complement of a fixed saddle
neighbourhood has exponentially small weight. The argument is a
fixed-\(N\) Laplace expansion.

There is a signed, elementary bound. Write \(N=N_2N_3\) and
\(n_*=\min(N_2,N_3)\). Expand \(\log(12-\mathsf A_\theta)\), with
\(\mathsf A_\theta\) the nearest-neighbour adjacency matrix carrying the
flat phases. If \(a_n(w)\) counts length-\(n\) walks from one site to
its translate by \((w_2N_2,w_3N_3)\) on the covering lattice, then

\[\begin{aligned}
\Omega_N(\theta)
 &=N\sum_{n\ge1}\frac1{n12^n}\sum_{w\ne0}
 a_n(w)[1-\cos(w\cdot\theta)],\\
0\le\Omega_N(\theta)&\le
\frac{3N}{n_*}\,3^{-n_*}.
\end{aligned}\tag{17}\]

The bound uses \(\sum_wa_n(w)\le4^n\) and \(n\ge n_*\) for nonzero
winding. Straight winding paths make the inequality strict whenever at
least one adjoint holonomy is nontrivial. At fixed \(N\), (16) is an
order-one background-dependent term as \(t\to0\). Both coupling shifts
and every local curvature operator vanish on these configurations.
Subtracting P(\(\alpha\)) for \(U_\theta\) and for \(1\) would bound
their difference by \(Ct^\alpha N\), contradicting (16). So
P(\(\alpha\)) as literally stated fails on fixed periodic planes with
arbitrary flat holonomies. The failure is an artifact of \(t\to0\) at
fixed \(N\): under the refinement scaling, with \(\min_jN_j\ge c_0/t\)
(fixed physical size), the winding term is \(O(t^{-1}e^{-c/t})\) and
lies inside every P(\(\alpha\)) remainder. The simplest repair is to add
the clause \(\min_jN_j\ge c_0/t\) to P(\(\alpha\)). The exponential
smallness comes from the massive mid-plane modes only; the coarse
theory's own holonomy potential is a separate object, not exponentially
small (Fable referee, 2026-09-27).

For \(N_j=L_j/a\) at fixed positive physical lengths and
\(t=\lambda_3a\), (17) is beyond every power of \(t\). Its origin is a
winding determinant; assigning it a local operator dimension would lose
its dependence on the torus. This is the precise exception to the bulk
irrelevance statement of §6.

A sufficient formulation for a theorem is the following. Work in local
gauge charts on the small-field set, including the treatment of global
holonomies and boundary conditions. For every plane size require a
unique contributing small saddle, a Hessian bounded below by a fixed
\(cI>0\), and exponentially summable connected kernels for its classical
value \(F\) and normalized one-loop amplitude \(A\). Constants and decay
rates must be independent of plane area. The normalized integral,
including all bridge denominators, Haar factors and images, must obey
the following after separating a finite-volume winding functional
\(W_N(U,t)\) (absent in the infinite-plane bulk expansion):

\[\begin{aligned}
\mathcal D(U,t)&=c_{\rm vac}(t)N+F(U)/t+A(U)-A(1)+W_N(U,t)+E(U,t),\\
|E(U,t)-E(1,t)|&\le C t\sum_p|X_p|^2+CNe^{-c/t},\\
|F-F_2-F_3-F_4|&\le C\varepsilon^3\sum_p|X_p|^2,\\
|A-A(1)-A_2|&\le C\varepsilon\sum_p|X_p|^2,
\qquad \varepsilon=\sup_p|X_p|\le t^{1/2-\delta}.
\end{aligned}\tag{18}\]

An \(E(1,t)\) constant is absorbed into \(c_{\rm vac}\). The quadratic
Taylor kernel of \(A\) must have the bulk coefficients (5), with
exponentially summable second spatial moments. Here \(F,A\) denote the
bulk local kernels; \(W_N\) collects their winding corrections and has
\(W_N(U_\theta,t)=\Omega_N(\theta)+O_N(t)\) after vacuum subtraction.
One further needs \(|W_N|\le C e^{-\kappa n_*}\sum_p(1+|X_p|^2/t)\) with
fixed \(\kappa>0\) in the chosen small-field charts. These conditions
justify both the order-\(t\) coefficient extraction and the derivative
expansion. They also give, after retaining the terms in (4), a remainder
bounded by

\[C\left(\frac{\varepsilon^3}{t}+\varepsilon+t\right)
\sum_p|X_p|^2+CNe^{-c/t}.\]

For the weaker unperturbed P(\(\alpha\)) inequality, it suffices to
replace the strong remainder line of (18) by
\(|E|\le Ct\sum_p(1+|X_p|^2/t)\) and use
\(|F-F_2|\le C\varepsilon\sum_p|X_p|^2\) and
\(|A-A(1)|\le C\sum_p|X_p|^2\). Then \(\alpha=1/2-\delta>2\delta\) for
\(0<\delta<1/6\), exactly as in Proposition 7, provided \(W_N\) is
retained or \(e^{-\kappa n_*}\le Ct^\alpha\) in the allowed
volume/scaling regime. The same estimates for the perturbed actions
produced by preceding steps are the additional iteration obligation.

The mass bound in (3) establishes the free reference's locality. The
remaining estimate is (18) for the \textbf{full normalized interacting
integral}, uniformly in the plane size and across the required gauge
charts, with \(W_N\) defined as the torus kernel minus the periodized
bulk kernels and the saddle's uniqueness supplied by
\(H\ge(4-C\varepsilon)I\) on the small-field set. Fixed-spin image
bounds and a finite-cube Laplace expansion supply local ingredients
only; the uniform estimate is the open item. The remainder exponents,
\(\alpha=\frac32-3\delta\) (strong) and \(\frac12-\delta\) (weak) with
\(\delta<\frac16\), suffice for one step.

\hypertarget{b.-the-coefficients-for-sun-and-su3}{%
\subsection{\texorpdfstring{7b. The coefficients for \(SU(N)\) and
\(SU(3)\)}{7b. The coefficients for SU(N) and SU(3)}}\label{b.-the-coefficients-for-sun-and-su3}}

\textbf{Result (Claude, 2026-09-29; written derivation, unrefereed).}
With the metric \(\langle X,Y\rangle=-2\operatorname{tr}XY\) used
throughout, the order-\(t\) coefficients (1) for \(G=SU(N)\) are those
of \(SU(2)\) multiplied by \(N/2\):
\[\delta_t^{12}=\delta_t^{13}=\frac{Nt}{48}\int_{\mathcal B}\frac qR+O(t^2),\qquad
\delta_t^{23}=\frac N2\,t\Bigl(-\frac1{12}+2\mathcal I\Bigr)+O(t^2).\]
For \(SU(3)\):
\(\delta_t^{12}=\delta_t^{13}=\frac t{16}\int_{\mathcal B}q/R>0\) and
\(\delta_t^{23}=t\bigl(-\frac18+3\mathcal I\bigr)<0\). The integrals
\(\int q/R\) and \(\mathcal I\) of (2) are group-independent, so the
signs of §§2--4 carry over.

\emph{Proof.} The coefficients are quadratic in the background and
one-loop, and every group factor in §§1--4 enters through the adjoint
action of the background on the fluctuations. Three facts make the
count.

\begin{enumerate}
\def\labelenumi{(\alph{enumi})}
\item
  For \(su(N)\) the Killing form is \(2N\) times the trace form, so
  \(\operatorname{tr}_{\rm adj}(\operatorname{ad}X)^2=2N\operatorname{tr}X^2=-N|X|^2\);
  for \(SU(2)\) this is the \(-2|X|^2\) used in (8). The cut shift (8)
  is linear in \(\operatorname{tr}_{\rm colour}(\operatorname{ad}X)^2\),
  and (6) and the exact bridge Gaussian (7) hold for any compact group
  (the zero-image factors
  \(\prod_{\alpha>0}\frac{\alpha(\xi)/2}{\sin(\alpha(\xi)/2)}\) of the
  heat kernel cancel the Haar Jacobian), so \(\delta^{12},\delta^{13}\)
  scale by \(N/2\).
\item
  The mid-face amplitude uses the Haar Jacobian
  \(j(Y)=\prod_{\alpha>0}[\sin(\alpha(Y)/2)/(\alpha(Y)/2)]^2\), whose
  logarithm is \(-\frac1{12}\sum_{\alpha>0}\alpha(Y)^2+O(Y^4)\), and
  \(\sum_{\alpha>0}\alpha(Y)^2=-\frac12\operatorname{tr}_{\rm adj}(\operatorname{ad}Y)^2=\frac N2|Y|^2\).
  For \(SU(2)\) the sum is \(|Y|^2\), which gave \(-t/12\); for
  \(SU(N)\) it gives \(-\frac N2\cdot\frac t{12}\).
\item
  The transverse coefficient is an \(\operatorname{Ad}\)-invariant
  quadratic form on the simple algebra \(su(N)\), hence a multiple of
  the Killing form, and it may be computed on a constant Cartan
  background \(Y\). There (10) splits by root spaces: \(Q_Y\) acts on
  the root space of \(\alpha\) as \(h_{\alpha(Y)}\), commutators of
  different root spaces are orthogonal to \(Y\), and the Cartan
  directions keep the free Hessian. Each root \(SU(2)\) is embedded
  isometrically, with \(T_3^{(\alpha)}=\frac i2(e_{ii}-e_{jj})\) of unit
  norm and charge 1 on its root space, so each positive root contributes
  one complex charged field with the Hessian (11) at \(B=\alpha(Y)\).
  The \(B^2\) coefficient of \(\log\det(8+K)\) therefore adds up to
  \(\sum_{\alpha>0}\alpha(Y)^2=\frac N2|Y|^2\) times the \(SU(2)\) value
  per unit \(B^2\). With (b), \(\delta^{23}\) scales by \(N/2\).
  \(\square\)
\end{enumerate}

For a general compact simple group the factor is the ratio of adjoint
Casimirs in the same metric, the one-loop pattern that also gives
\(b_0=11N/(48\pi^2)\) in the
\href{four-dimensional-parallel-log.md}{four-dimensional note}. The
scaling covers the bulk coefficients only. The dimension-six content of
\(SU(3)\) adds the \(d_{abc}\) cubic noted in §6, and the finite-torus
winding term (16) depends on flat holonomies and \(\mathbb Z_3\)
sectors, which the \href{su2-midplane-small-field.md}{small-field note},
§16, treats separately.

\hypertarget{consequence-for-state}{%
\subsection{8. Consequence for STATE}\label{consequence-for-state}}

Atlas cell 2 now has a formal full-plane relative-order-\(t\)
calculation: the cut couplings shift positively, the transverse coupling
negatively, and all other bulk smooth-field terms start at dimension
six. The winding term (16) obstructs P(\(\alpha\)) as literally stated
at fixed plane size \(N\); in the refinement scaling it lies in the
remainder, and P(\(\alpha\)) gains the clause \(\min_jN_j\ge c_0/t\).
The next obligation for this cell is the normalized estimate (18), with
winding control, followed by stability under iteration. The infrared
spectral-gap obligation remains separate. For \(SU(3)\) the bulk
coefficients are the \(SU(2)\) ones times \(\frac32\) (§7b). STATE
continues to route work through the atlas; the coefficients and their
derivation live here.
\end{document}
